\section{Separating two subfield ends}
The following is the complete previously proved intersection argument,
now placed next to the common trace obstruction it uses. Its earlier
independent acceptance and finite diagnostic are provenance, not an
unpriced substitute for the proof below.

\begin{lemma}\label{E.TranslatedSubfields}
Let $F_1,F_2$ be subfields of a field $K$, with $F_1\cap F_2=F$.
If $t\in K\setminus(F_1+F_2)$, then
\[
 L(F_1)\cap U(t)L(F_2)U(-t)
 =\{[U(b)]:b\in F\},\qquad U(t)=\begin{pmatrix}1&t\\0&1\end{pmatrix},
\]
in every characteristic.
\end{lemma}
\begin{proof}
Write a determinant-one representative in the intersection as
$M=\left(\begin{smallmatrix}a&b\\c&d\end{smallmatrix}\right)$ over $F_1$.
Its representative over the other field differs only by a scalar $\pm1$,
which can be absorbed. Thus
\[
 U(-t)MU(t)=\begin{pmatrix}
 a-tc&b+t(a-d)-t^2c\\c&d+tc
 \end{pmatrix}
\]
has entries in $F_2$. Hence $c\in F$. If $c\ne0$, the first entry
would put $t$ in $F_1+F_2$, so $c=0$. Now $a,d\in F$; if $a\ne d$,
the upper-right entry gives the same contradiction. Thus $a=d=\pm1$
and $b\in F$. Conversely every displayed unipotent belongs to both
groups and commutes with $U(t)$.
\end{proof}

\begin{theorem}\label{E.CoprimeConstruction}
Let $p$ be odd and $e_1,e_2>1$ coprime odd integers. Put
$F_i=F_{p^{e_i}}$, $K=F_{p^{e_1e_2}}$. Choose primitive $j_i\in F_i$
by Theorem~\ref{E.PrimitiveSquare}, put $u=j_1+j_1^{-1}$,
$v=j_2+j_2^{-1}$ and choose $D\ne0$ with
$D^2=u^2v^2-4(u^2+v^2)$. The projective matrices
\[
 \alpha=\frac12\begin{pmatrix}u&u+2\\u-2&u\end{pmatrix},\quad
 \beta=\frac12\begin{pmatrix}2+u&u\\-u&2-u\end{pmatrix},\quad
 \gamma=\frac1{2u}\begin{pmatrix}v(u-2)&D\\D&v(u+2)\end{pmatrix}
\]
give a chiral rank-four polytope with full group $L(K)$ and type
$((p^{e_1}-1)/2,p,(p^{e_2}-1)/2)$.
\end{theorem}
\begin{proof}
The odd-degree argument used above shows that $u^2$ and $v^2$ generate
their full respective fields over $F_p$; in particular $u,v\ne0$ and
$u^2,v^2\ne4$. The matrices have determinant one, with traces $u,2,v$.
Their eigenvalues give the first and last orders, while $\beta=I+N$
with $N\ne0$, $N^2=0$, gives order $p$. These orders imply the adjacent
cyclic intersections are trivial. Direct multiplication gives
$\alpha\beta=J$ and zero traces for $\beta\gamma,\alpha\beta\gamma$,
so all interval-involution relations hold.

We prove the remaining intersection without assuming generation.
If $x\in F_1$, $y\in F_2$ and $xy\in F_1+F_2$, then $x\in F_p$
or $y\in F_p$. Indeed, apply the $p^{e_1}$-Frobenius $\tau$ to
$xy=a+b$, $a\in F_1,b\in F_2$, and subtract. If $y^\tau\ne y$,
then $x=(b^\tau-b)/(y^\tau-y)\in F_2$, hence in $F_p$; otherwise
$y\in F_p$. Put $z=D/(uv)$. It satisfies
$z^2=1-4/u^2-4/v^2$ and does not belong to $F_1+F_2$.
For if $z=x+y$ with $x\in F_1,y\in F_2$, the square equation puts
$2xy$ in that sum. The preceding observation forces $z$ into one of
the two fields, and then puts $u^2$ or $v^2$ in $F_p$, a contradiction.

With $M^T=T^{-1}MT$, put
\[
 T=\begin{pmatrix}1&2/u\\-1&0\end{pmatrix},\qquad
 A(w)=\begin{pmatrix}1&(2-w)/w\\-w&w-1\end{pmatrix},\qquad
 t=\frac D{2uv}+\frac1v-\frac12-\frac1u.
\]
Direct multiplication gives
\[
 \beta^T=U(1),\quad\alpha^T=A(u),\qquad(\gamma^T)^{U(t)}=A(v),
\]\[
 \gamma^T=\begin{pmatrix}
 (-D+v(u+2))/(2u)&-D/u^2\\-v&(D+v(u-2))/(2u)
 \end{pmatrix}.
\]
The last identity also follows from its upper-left entry $1$, lower-left
entry $-v$, trace $v$ and determinant one. Since $t\notin F_1+F_2$,
Lemma~\ref{E.TranslatedSubfields} puts the two actual end groups'
intersection inside $\langle U(1)\rangle$. Both contain that group,
so their intersection is exactly $\langle\beta\rangle$.

For generation use the retained maximal-subgroup list. The unipotent
$\beta^T$ fixes only infinity, whereas $A(u)$ does not fix infinity,
so the group is in no point stabilizer. The dihedral maximal orders
$|K|\pm1$ are prime to $p$ and cannot contain $\beta$. The exceptional
$A_4,S_4,A_5$ have no element of order at least $13$, whereas
$|\alpha|=(p^{e_1}-1)/2\ge13$. Finally the conjugacy invariant
$\operatorname{tr}(M)^2/\det(M)$ takes the values $u^2,v^2$ on
$\alpha,\gamma$. A conjugate subfield group containing them must
therefore contain both $F_1,F_2$ in its field of definition, and hence
all of $K$. No proper maximal subgroup remains.

A reversing automorphism is semilinear by the retained automorphism
theorem. Its field part fixes $u^2$ and $v^2$, so fixes $K$ and is
trivial. The trace obstruction in Theorem~\ref{E.TraceExtension}, or
its final two-line centralizer calculation alone, now applies:
\[
 \operatorname{tr}(\alpha\gamma)-uv/2=D/2\ne0.
\]
Thus no projectivity can simultaneously invert $\alpha$ and fix
$\gamma$. There is no reversal, so the full rotation correspondence
gives the asserted chiral polytope.
\end{proof}

Every odd integer $d>1$ with at least two distinct prime divisors admits
a factorization $d=e_1e_2$ into coprime odd integers greater than one,
by separating its prime-power factors. This theorem supplies all such
exponents. Together with the earlier exclusions, the nonprime-field
classification is consequently: $L(p^d)$, $d>1$, occurs in rank four
exactly for even $q>4$, for $q\equiv1\pmod4$ with $q>9$, or for
$q\equiv3\pmod4$ when $d$ is not a prime power. This reproduces the
previously accepted Section~6 completion. It is kept separate from
the source's established principal theorem and earns no new question
or historical-priority claim in this edition.
